% Source: https://artofproblemsolving.com/community/c3h1568113_mathcounts_questions
% Stable identity: aops-topic-1568113-q05
Whether $n$ is in $K$ depends only on the remainder of $n$ modulo $60=\operatorname{lcm}(3,4,5)$, so $K$ repeats with period $60$.

\textbf{Count in one period $1,\dots,60$.} Numbers divisible by neither $3$ nor $4$: $60-20-15+5=30$. Numbers divisible by $3$ or $4$ and also by $5$, i.e.\ by $15$ or $20$: $4+3-1=6$. So each block of $60$ consecutive integers contains $36$ members of $K$.

\textbf{Locate the $79$th.} $79=2\cdot36+7$, so it is the $7$th member of $K$ in the block $121,\dots,180$, i.e.\ $120$ plus the $7$th member of $K$ in $1,\dots,60$. Those members begin
\[
1,\ 2,\ 5,\ 7,\ 10,\ 11,\ 13,\ \dots
\]
so the $7$th is $13$, and the $79$th member of $K$ is $120+13=\boxed{133}$.
